\chapter{绪论}

\section{研究背景与意义}
\subsection{新型内存介质与持久内存的发展}
说明持久内存、CXL 扩展内存等介质对容量、持久性和访问路径的改变，以及这些变化对数据库数据组织方式提出的新要求。
\subsection{远程内存与解耦内存架构}
介绍 RDMA、CXL 内存池和计算—内存解耦架构，分析远程访问时延、带宽和内存节点计算能力约束。
\subsection{云原生数据库中的计算—内存—存储解耦}
说明计算节点（CN）、内存节点（MN）和存储节点（SN）独立扩展的趋势，以及资源非对称性带来的管理功能放置问题。
\subsection{新型内存体系对数据管理机制的挑战}
将全文问题归纳为三类失配：动态数据位置与显式映射元数据之间的\textbf{数据表示失配}，本地数据结构与远程访问成本之间的\textbf{数据布局失配}，以及单体管理组件与异构资源之间的\textbf{管理功能失配}。

\section{国内外研究现状}
\subsection{持久内存数据管理}
综述原地更新、日志、写时复制、地址映射与崩溃一致性机制。
\subsection{远程与解耦内存数据组织}
综述 RDMA 数据结构、远程指针追踪、批量访问和内存池化研究。
\subsection{多版本数据管理}
综述链式、连续式和分层式版本布局，以及版本索引与垃圾回收机制。
\subsection{云原生数据库缓存管理}
综述 MN 中心、CN 中心及分布式缓存管理方案。
\subsection{现有研究的不足}
归纳传统方案对本地、均衡、紧耦合成本模型的依赖，并明确本文的研究边界。

\section{本文研究问题}
\begin{enumerate}
  \item \textbf{RQ1：}如何降低动态数据位置变化带来的数据寻址和元数据管理开销？
  \item \textbf{RQ2：}如何设计适应远程访问时延、带宽及版本分布变化的数据组织机制？
  \item \textbf{RQ3：}如何根据计算、内存和存储资源的异构属性重新划分数据管理功能？
\end{enumerate}

\section{本文研究内容}
\subsection{硬件辅助的数据寻址机制}
对应 ReCoW：复用虚拟—物理地址映射能力，降低写时复制的数据位置跟踪成本。
\subsection{网络感知的自适应多版本数据组织机制}
对应 MiTDM 及其博士论文扩展：采用层次化版本块并研究网络与负载驱动的粒度调节。
\subsection{资源感知的分离式缓存管理机制}
对应 Duet：按资源属性拆分替换策略、换入和写回任务。

\section{技术路线与论文结构}
\subsection{技术路线}
按照“发现失配—建立成本模型—重构机制—原型实现—实验验证”的路线，研究范围由单节点持久内存扩展到计算—内存解耦，再扩展到计算—内存—存储解耦。
\subsection{论文结构}
概述第二章至第七章的内容及逻辑关系，并给出“数据在哪里—数据如何组织—数据由谁管理”的主线图。

\section{本文主要创新点}
\begin{enumerate}
  \item 提出基于虚拟地址重映射的硬件辅助数据寻址机制，降低动态数据位置管理的软件元数据开销。
  \item 提出面向解耦内存的层次化多版本数据组织机制，并进一步研究网络条件和版本分布驱动的自适应粒度选择。
  \item 提出面向异构资源的分离式缓存管理架构，将计算密集、时延敏感和带宽密集任务放置到更匹配的资源层。
\end{enumerate}

\section{本章小结}
总结研究动机、核心问题、技术路线和全文结构。
